\section{Problem Formulation and Threat Model}
\label{sec:problem}

\subsection{Objects, Scores, and Decisions}

Let $x\in\mathbb{R}^{d}$ denote a static PE feature vector with $d=2381$ in the
available pilot. A static classifier $f_{\mathrm{s}}$ produces a bounded score
$\pstatic=f_{\mathrm{s}}(x)\in[0,1]$. The implementation uses this score as a
probability-like output, but the fusion design does not assume perfect
calibration. If dynamic analysis is invoked, an independently specified mapping
$f_{\mathrm{d}}$ converts observed behavior $z$ into a bounded dynamic score
$\pdynamic=f_{\mathrm{d}}(z)\in[0,1]$. The meaning of $\pdynamic$---calibrated
probability or risk score---must be fixed before evaluation.

The architecture contains a conceptual uncertainty gate $g(\pstatic)$. A common
deployment form uses two validation-selected thresholds $\tau_{\mathrm{low}}$ and
$\tau_{\mathrm{high}}$: scores outside the interval can receive an early static
decision, while scores inside it are escalated for behavioral observation. The
present measured artifacts do not evaluate this gate or its coverage. In
particular, they do not establish how often a confident static error would bypass
dynamic analysis. The fusion protocol instead evaluates all supplied, aligned
static--dynamic pairs; any future gate study must report both conditional quality
and the fraction of traffic escalated.

For each evaluated score $p$, a threshold $t$ maps the score to
$\widehat{y}=\mathbb{1}[p\geq t]$. Thresholds are selected on the validation
partition for a target true-positive rate (TPR) of 0.95 and are then frozen for
test evaluation. The test TPR is still reported because a validation target does
not ensure exactly the same operating point on test. With true positives (TP),
false positives (FP), true negatives (TN), and false negatives (FN),
\begin{equation}
\mathrm{TPR}=\frac{\mathrm{TP}}{\mathrm{TP}+\mathrm{FN}},\qquad
\mathrm{FPR}=\frac{\mathrm{FP}}{\mathrm{FP}+\mathrm{TN}}.
\label{eq:tpr-fpr}
\end{equation}
ROC--AUC and average precision (AP) summarize ranking, while FPR and TPR describe
the selected operating point \cite{fawcett2006roc,saito2015precision}.

\subsection{Active-Learning Setting}

Let $\mathcal{L}_r$ and $\mathcal{U}_r$ be the labeled and unlabeled training
pools at round $r$. The model $f_r$ is fitted only on $\mathcal{L}_r$. An
acquisition function $Q_r(x)$ ranks $x\in\mathcal{U}_r$, and the top $b$ rows are
sent to an oracle. In the pilot, the oracle is simulated by revealing an existing
label only after selection. This permits controlled comparison of query policies;
it does not measure analyst disagreement, annotation time, or monetary cost.

All strategies within a seed share the same initial labeled set, training pool,
validation set, test set, and cumulative label budget. The test partition is not
used for querying, parameter selection, threshold selection, or stopping. The
primary estimand for RQ1 is the paired test-metric difference between acquisition
strategies at the same seed and budget. Because the five runs share one test set,
seed variability describes sensitivity to initialization and queried samples, not
a population confidence interval over independently drawn test sets.

\subsection{Explanation Object and Intervention}

For a fitted tree model and input $x_i$, let
$\boldsymbol{\phi}_i=(\phi_{i1},\ldots,\phi_{id})$ be the signed TreeSHAP
contribution vector, excluding the bias term. Native TreeSHAP operates on the raw
log-odds scale and must satisfy the additive reconstruction
\begin{equation}
f_{\mathrm{raw}}(x_i)=\phi_{i0}+\sum_{j=1}^{d}\phi_{ij}
\label{eq:additivity}
\end{equation}
within numerical tolerance. This identity verifies the computation, not the
semantic correctness of an explanation.

To test model sensitivity, a mask $M_i^{(k)}$ selects $k$ coordinates for sample
$i$. The intervention replaces each selected value by the corresponding
arithmetic mean $\mu_j$ estimated from the training partition:
\begin{equation}
\widetilde{x}_{ij}=\begin{cases}
\mu_j, & j\in M_i^{(k)},\\
x_{ij}, & \text{otherwise}.
\end{cases}
\label{eq:mask}
\end{equation}
The top-SHAP mask contains the $k$ largest $|\phi_{ij}|$ values for each sample.
The control selects $k$ distinct coordinates uniformly at random. This is a
feature-space deletion test. It does not claim that $\widetilde{x}_i$ is produced
by a valid PE or a semantics-preserving binary edit.

\subsection{Threat Model and Scope}

The motivating adversary distributes Windows PE malware and may employ packing,
obfuscation, benign-looking metadata, delayed execution, intermittent activity,
or environment checks. Such behavior can weaken either evidence source: static
features may hide functionality, while a sandbox may not observe the payload
within its window. The current design does not claim robustness to a fully
adaptive attacker with white-box access to the static model, query function, and
fusion rule. Adversarial feature manipulation, explanation manipulation, and
active-learning poisoning remain outside the measured scope.

The defender is assumed to have a static feature extractor, a LightGBM model, an
optional isolated execution environment, and an oracle capable of assigning
labels to selected samples. A valid PID evaluation additionally assumes that
static and dynamic scores are aligned to the same sample and observation session,
that sequence order is meaningful, and that validation and test sessions do not
share controller state. These assumptions are enforced structurally by the
runner but cannot be satisfied by inventing $\pdynamic$ from the class label,
perturbing $\pstatic$ with random noise, or replaying an untraceable demonstration
log.

Finally, the empirical scope is narrower than the motivating threat model. The
audited CSV has binary labels but no verified original file identities, family
metadata, timestamps, or ransomware-specific label definition. We therefore use
``positive'' and ``negative'' when discussing measured classes. Ransomware is the
system motivation and intended deployment domain, not a provenance claim about
every evaluated row.
